Mit diesem Projekt liefert das fiktive Ingenieurbüro „B·E·Sondere Lösung UG`` 
für den fiktiven Kunden „Pay With Fun GmbH`` einen Prototypen-Konzept für einen Coasterbot.
Dabei ist das explizite Ziel Grundfunktionen in dem Prototypen umzusetzen.
Diese Grundfunktionen umfassen folgendes:
\begin{itemize}
    \item Navigation auf einem Tischoberfläche
    \item Ausgabe von Getränkeuntersetzer
    \item Lokalisierung von Getränkeuntersetzer
    \item Aufnahme von Getränkeuntersetzer
\end{itemize}
Zur Umsetzung dieser Grundfunktionen dürfen bereitsbestehende Konzepte zu mobilen Robotern verwendet werden.
Innerhalb des Projektes sind weitere Ausbaustufen auszuarbeiten und mit einem Testkonzept zu versehen.

Zu diesem Prototypen sind folgende Liefergegenstände bereitzustellen:
\begin{itemize}
    \item Eine Aufbauanleitung
    \item Eine Stückliste (Bill of Materials, BOM)
    \item Die zur Fertigung notwendigen Modelle als stl-Dateien
    \item Ein Simulationsmodell zur Weiterentwicklung
    \item Der Sourcecode zum Prototypen
\end{itemize}
Mit dem Upload dieser Liefergegenstände über das Dokumenten-Management-System (DMS) 
der Wilhelm-Büchner-Hochschule (WBH) gelten sie als ausgeliefert. 
Ein physischer Aufbau wird nicht separat geliefert werden. 

Nicht als Bestandteil des Projekts gelten:
\begin{itemize}
    \item Ein Serienfertiges Endprodukt
    \item EMV Betrachtungen
    \item Intelligentes Schwarmverhalten
    \item Funktionen außerhalb der Grundfunktionen
    \item Konzeption einer Vermarktung des Produkts
\end{itemize}

Damit wird das Projekt als „PA2 - Machbarkeitsprojekt`` im Maßstab der Wilhelm-Büchner-Hochschule geführt.
Der Projekstart erfolgte am 21.07.2026 und das Projektende ist am 07.11.2026 mit der Abgabe dieses Endbericht. 
Damit beträgt der angesetzte zeitliche Rahmen 3,5 Monate angesetzt. 
